\section*{الفصل \ref{ch:b2:plant-plasticity} --- تكيفات النبات واللدونة النمطية}

\begin{solution}{exo:b2:plant-plasticity:1}
اللدونة: \omterm{def:b2:meiosis-heredity:vocabulary}{نمط وراثي} واحد، وأنماط ظاهرية عدة بحسب
المحيط (أوراق شمس وظل على بلوطة واحدة). \omterm{def:b2:plant-plasticity:plasticity}{ومعيار الاستجابة}:
منحني \omterm{def:b2:meiosis-heredity:vocabulary}{النمط الظاهري} \omterm{def:b2:meiosis-heredity:vocabulary}{لنمط وراثي} بدلالة متغير
بيئي (ثخن الورقة بدلالة الضوء). \omterm{def:b2:plant-plasticity:plasticity}{والتأقلم}: تعديل
فيزيولوجي عكوس (تقسية الجاودار على الصقيع في الخريف).
\omterm{def:b2:plant-plasticity:plasticity}{والتكيف}: ملاءمة موروثة منتخبة على مدى الأجيال (استقلاب
CAM عند صبار).
\end{solution}

\begin{solution}{exo:b2:plant-plasticity:2}
غمد سويقة سليم مضاء من جانب: ينحني نحو الضوء. والقمة
مغطاة بغطاء معتم: لا انحناء --- فالقمة تدرك
الضوء. والقمة مغطاة بغطاء شفاف: ينحني --- فالغطاء نفسه
غير ضار. والقاعدة محجوبة بأنبوب معتم والقمة مكشوفة: ينحني ---
فالإدراك في القمة والاستجابة تحتها.
\end{solution}

\begin{solution}{exo:b2:plant-plasticity:3}
\omterm{prop:b2:plant-plasticity:sunshade}{ورقة الشمس}: صغيرة سميكة، بطبقتين أو ثلاث عمادية، وبشرة
سميكة، وثغور كثيرة، وتركيب ضوئي أقصى عالٍ ونقطة تعويض
عالية. \omterm{prop:b2:plant-plasticity:sunshade}{وورقة الظل}: كبيرة رقيقة، بطبقة عمادية واحدة،
وكلوروفيل أكثر لكل غرام، وتنفس منخفض، ونقطة تعويض منخفضة.
ويُتخذ الاختيار والورقة بدائية في البرعم، من
الضوء الذي تتلقاه؛ والورقة الناضجة لا تستطيع التغير.
\end{solution}

\begin{solution}{exo:b2:plant-plasticity:4}
تأخذ الخليتان الحارستان البوتاسيوم والماء، فتنتفخان، ولمّا كان
جداراهما الداخليان أثخن تقوستا متباعدتين ففتحتا المسمّ. والضوء الأزرق
و$\mathrm{CO_2}$ الداخلي المنخفض يفتحانهما في الصباح؛ \omterm{def:b2:plant-flowering:dormancy}{وحمض
الأبسيسيك}، من جذور في تربة تجف أو من الورقة، يغلقهما في
الجفاف، ويغلقهما كذلك هواء شديد الجفاف.
\end{solution}

\begin{solution}{exo:b2:plant-plasticity:5}
$\varphi = R/(R + 0.77\,FR)$ مع $FR = 1$: فالقيمة $1.15$ تعطي $0.60$؛ و$0.7$ تعطي
$0.48$؛ و$0.2$ تعطي $0.21$؛ و$0.05$ تعطي $0.06$. والعتبة $0.4$ تقع
بين الفرجة ($0.48$) والمظلة ($0.21$): فيشتعل \omterm{thm:b2:plant-plasticity:phytochrome}{تجنب الظل}
تحت المظلة.
\end{solution}

\begin{solution}{exo:b2:plant-plasticity:6}
$E = 0.25\times 0.025 = \qty{6.25}{mmol.m^{-2}.s^{-1}}$؛ و$A =
0.25\times 100/1.6 = \qty{15.6}{\micro mol.m^{-2}.s^{-1}}$؛ و$A/E =
2.5\times 10^{-3}$ مول كربون لكل مول ماء، أي 400 جزيء
ماء لكل كربون، أي $400\times 18/12 = \qty{600}{g}$ من الماء لكل
غرام كربون. وفي 12 ساعة: $6.25\times 10^{-3}\times 43200 =
\qty{270}{mol/m^2}$، أي \qty{4.9}{L} في المتر المربع.
\end{solution}

\begin{solution}{exo:b2:plant-plasticity:7}
$E = 0.25\times 0.005 = \qty{1.25}{mmol.m^{-2}.s^{-1}}$، و$A$
بلا تغير: فتكون $A/E = 0.0125$، أي 80 جزيء ماء لكل كربون، أي \qty{120}{g/g} ---
أي خمس كلفة النهار.
\end{solution}

\begin{solution}{exo:b2:plant-plasticity:8}
$v = 2r^{2}\Delta\rho g/9\eta = 2\times 4\times 10^{-12}\times
400\times 9.81/0.09 = \qty{3.5e-7}{m/s}$: أي \qty{0.35}{\micro m/s}،
ونحو \qty{30}{s} لعبور الخلية. وعلى الكلينوستات يتغير اتجاه
الجاذبية بالنسبة إلى الخلية كل بضع ثوان من الترسيب،
فلا تتراكم الحبيبات على جانب واحد أبدًا ويكون المنبه
المتوسط على مدى زمن مكاملة النبات (دقائق) صفرًا.
\end{solution}

\begin{solution}{exo:b2:plant-plasticity:9}
متوسط الاستطالة في ساعتين $0.04$؛ والجانبان $1.2\times 0.04 = 0.048$
و$0.8\times 0.04 = 0.032$؛ ومنه $\theta = 20\times 0.016/2 = 0.16$
راديان، أي نحو \qty{9}{\degree}. ويتراكم الانحناء بمقدار $0.08$ راديان في
الساعة؛ فيستغرق $\pi/2$ نحو \qty{20}{h}.
\end{solution}

\begin{solution}{exo:b2:plant-plasticity:10}
المدخرات تتيح \qty{50}{cm} من الساق. وفي حقل القراص: \qty{30}{cm} في
7.5 يوم مقابل \qty{30}{mg} --- فتبلغ البادرة الضوء ولها
مدخرات فائضة. وفي الغابة: \qty{10}{m} خارج المتناول؛ فتنفق البادرة
\qty{50}{mg} على نصف متر من ساق شاحبة في اثني عشر
يومًا وتموت مصفرّة، في حين أن استراتيجية احتمال الظل --- بضع
أوراق رقيقة، ونمو بطيء --- ربما أبقتها حية سنوات.
\end{solution}

\begin{solution}{exo:b2:plant-plasticity:11}
يتكون الثلج أولًا في الفراغات خارج الخلوية، حيث يكون المحلول
أرق؛ وضغط البخار فوق الثلج أدنى منه فوق ماء
الخلية، فيغادر الماء الخلية لينضم إلى الثلج فتتجفف الخلية
وتنهار أغشيتها. وأسبوع من البرد يحمّل
الخلية سكريات وبروتينات واقية تخفض نقطة تجمدها
وتثبّت الأغشية، ويصنع بروتينات مانعة للتجمد
تحد نمو البلورات. أما إبقاء البرنامج مشتعلًا طوال السنة فيحوّل
الكربون من النمو إلى الحماية ويبطئ الانقسام: فالنبات المقسّى
دائمًا نبات صغير.
\end{solution}

\begin{solution}{exo:b2:plant-plasticity:12}
إشارة الأحمر البعيد صادقة حين يمكن أن يُعلى على الجار
وكاذبة حين يكون شجرة؛ والأسبوع الدافئ صادق في نيسان و
كاذب في شباط؛ وفي محصول كثيف يكون جيران كل نبات
أنداده، فتكون الإشارة صادقة شكلًا لكن عديمة النفع، وينتخب المربّون
نباتات تتجاهلها. \omterm{def:b2:plant-plasticity:plasticity}{ومعيار الاستجابة} سجل
كم مرة قال كل إشارة الحقيقة في ماضي السلالة، وهو
يخفق كلما اختلف الحاضر عن ذلك الماضي.
\end{solution}

\begin{solution}{pb:b2:plant-plasticity:1}
\textbf{1.} المكشوف: $1.15/(1.15 + 0.77) = 0.60$؛ والمظلة: $0.2/(0.2 +
0.77) = 0.21$.
\textbf{2.} $\varphi$ منخفضة: أي نباتات فوقها --- فتستطيل. أما تحت
قماش تظليل فتبقى $\varphi$ عند $0.60$: أي يوم غائم، لا جار؛ فلا
تجنب ظل، بل نمو أبطأ لقلة الضوء فقط.
\textbf{3.} $1 + 2(0.60 - 0.21) = 1.78$: أي \qty{2.7}{cm} في اليوم.
\textbf{4.} \qty{37}{cm} في مقابل \qty{21}{cm}.
\textbf{5.} انحناء نحو الفرجة (\omterm{prop:b2:plant-plasticity:phototropism}{انتحاء ضوئي}، عبر مستقبل الضوء
الأزرق الفوتوتروبين) واستطالة أسرع وإعادة توجيه الأوراق
في الجانب الأعلى $R{:}FR$ (عبر الفيتوكروم).
\textbf{6.} أوراق الظل مبنية كبيرة رقيقة لالتقاط الضوء
الشحيح بثمن بخس؛ وفي الشمس الكاملة تفرط في الحرارة وتُثبَّط ضوئيًا وتحترق،
وعلى النبات أن يعوّضها بأوراق شمس.
\textbf{7.} \qty{2.7}{cm} في اليوم على منطقة \qty{15}{mm} تساوي
\qty{0.11}{cm} في الساعة، أي استطالة نسبية $0.074$ في الساعة؛
فالجانب المظلل $1.3\times 0.074 = 0.096$، والمضاء $0.7\times 0.074 =
0.052$.
\textbf{8.} $\theta = 15\times(0.096 - 0.052)/2 = 0.33$ راديان، أي نحو
\qty{19}{\degree} في ساعة.
\textbf{9.} القيمة \qty{60}{\degree} هي $1.05$ راديان: أي نحو ثلاث ساعات.
\textbf{10.} $v = 2\times 6.25\times 10^{-12}\times 400\times
9.81/0.09 = \qty{5.5e-7}{m/s}$، أي \qty{0.55}{\micro m/s}: و\qty{22}{s}
لعبور الخلية.
\textbf{11.} \omterm{prop:b2:plant-plasticity:phototropism}{الانتحاء الضوئي}، \omterm{thm:b2:plant-plasticity:phytochrome}{وتجنب الظل} نفسه يرفع
الزاوية التي يرضى الفرع بالنمو عندها: فالضوء هو المورد
المتنازع عليه، والساق التي تستقيم ضد الجاذبية كانت
لتعود إلى الظل.
\textbf{12.} ينحني الجذر نازلًا حتى يصير شاقوليًا: إذ تستقر حصى التوازن
إلى الجانب السفلي، ويتراكم الأوكسين هناك، وفي الجذر
يثبّط الأوكسين الزائد النمو، فينمو الجانب العلوي أسرع. فإعادة
التوزيع نفسها، وحساسية الخلايا معاكسة.
\textbf{13.} $E = 0.2\times 0.015 = \qty{3}{mmol.m^{-2}.s^{-1}}$؛ وعلى
$2\times 10^{-3}$ \unit{m^2} و\qty{50400}{s}: أي \qty{0.30}{mol}،
أي \qty{5.4}{g} من الماء في اليوم.
\textbf{14.} $A = 0.2\times 80/1.6 = \qty{10}{\micro mol.m^{-2}.s^{-1}}$؛
وفي اليوم $0.50$ مول كربون في المتر المربع؛ و$A/E = 3.3\times
10^{-3}$: أي 300 جزيء ماء لكل كربون، أي \qty{450}{g/g}.
\textbf{15.} $0.50\times 2\times 10^{-3} = \qty{1}{mmol}$: أي \qty{12}{mg}
من الكربون في اليوم.
\textbf{16.} في الفرجة: $E = \qty{6}{mmol.m^{-2}.s^{-1}}$، أي \qty{10.9}{g} في
اليوم؛ والكفاءة \qty{900}{g/g}.
\textbf{17.} $E = 0.02\times 0.03 = \qty{0.6}{mmol.m^{-2}.s^{-1}}$
(أي عشرة أضعاف أدنى)؛ و$A = 0.02\times 200/1.6 = \qty{2.5}{\micro
mol.m^{-2}.s^{-1}}$ (أي أربعة أضعاف أدنى). فالماء يهبط أكثر: إذ
تسحب الورقة $\mathrm{CO_2}$ الداخلي فيتسع تدرج الكربون
بينما لا يتسع تدرج الماء.
\textbf{18.} \qty{4.25}{g} من الماء، منها \qty{0.85}{g} يمكن فقدها؛ وعند
\qty{10.9}{g} في اليوم (أي $0.78$ غ في الساعة) تذبل في نحو ساعة.
\textbf{19.} نمط CAM بطيء ويحتاج إلى نسيج تخزين عصاري؛ والبادرة
المتسابقة إلى الضوء تحت مظلة، حيث فقد الماء صغير
أصلًا، تحتاج إلى سرعة لا إلى اقتصاد.
\textbf{20.} $100/2 = \qty{50}{cm}$.
\textbf{21.} حقل القراص: \qty{40}{cm} في 15 يومًا مقابل \qty{80}{mg} ---
فتبلغ الضوء. والغابة: \qty{8}{m} لا يمكن بلوغها؛ فتنفق
مدخراتها على نصف متر وتموت. والأفضل في الغابة: أن تبقى قصيرة،
وتصنع أوراق ظل رقيقة وتنتظر فرجة --- وهي طريقة الزان.
\textbf{22.} البتولا: معيار حاد، تستطيل بقوة عند
$\varphi$ المنخفضة، إذ إن جيران الرائدة من قياسها. والزان:
معيار مسطح، باستطالة قليلة واحتمال كثير، إذ إن جيرانه
أشجار.
\textbf{23.} عطّل استجابة \omterm{thm:b2:plant-plasticity:phytochrome}{تجنب الظل} --- بإبقاء
الفيتوكروم B فعالًا أو بنزع \omterm{prop:b2:cell-differentiation:transcriptional}{عوامل الاستنساخ} التي
يكبحها --- وتحقق من موعد الإزهار والإنبات وقوة
الساق، فالمسلك نفسه يتحكم فيها.
\textbf{24.} النبات الذي لا يستطيع استشعار الجيران يُعلى عليه
كلما كان في وسعه أن يربح؛ والنبات الذي يصدّق الإشارة دائمًا
يبدد نفسه تحت كل شجرة. والمعيار الرابح يراهن على
الإشارة بقدر ما كانت، في ماضي السلالة،
محقة.
\textbf{25.} $\varphi = 0.60$ في المكشوف و$0.21$ تحت المظلة؛
و\qty{37}{cm} بعد 14 يومًا؛ و\qty{5.4}{g} من الماء في اليوم في الظل،
و\qty{10.9}{g} في الفرجة.
\end{solution}
